\documentclass[11pt,a4paper]{article}
\usepackage[margin=25mm]{geometry}
\usepackage{amsmath,amssymb,amsthm,booktabs}
\newtheorem{lemma}{Lemma}
\newtheorem{theorem}[lemma]{Theorem}
\newcommand{\gen}[1]{\langle#1\rangle}
\title{Enumerating a string from its two ends}
\author{Human Compression project}
\date{13 October 2026}
\begin{document}\maketitle

We retain the complete proofs of enumeration by successive stabilizers
and of its centralizer and double coset implementation acquired in the
L071 recovery. The present change is to assign labels in the order
$1,r,2,r-1,\ldots$ for a fixed rank $r$. This can expose commuting
restrictions before choosing adjacent marks. The argument works for any
fixed order of the labels. It does not assert that this choice is faster
on every group.

\paragraph{Premises.}
Let $G\lhd A$ be finite permutation groups. We assume that $G$ has been
identified as the specified nonabelian simple group and that conjugation
induces a surjection $A\longrightarrow\operatorname{Aut}(G)$. An
independently known value of $|\operatorname{Aut}(G)|$, together with
$|A:C_A(G)|$ of that value, verifies the latter assertion. A group order
by itself identifies neither $G$ nor its full automorphism action.
For the finite control below the identification of
\texttt{PrimitiveGroup(351,7)} with $G_2(3)$ is the retained standard
group-library identification. The full automorphism group is computed
from this group, and its inner automorphism subgroup is transported with
it into one faithful permutation action. Exact GAP algorithms for group
orders, intersections, conjugacy classes, automorphisms, double cosets,
Sylow subgroups, and subgroup classes remain software premises. We do
not replace those algorithms by the order checks in the driver.

The intersection property for marks $t_i$, $i\in I$, is
\[
 \gen{t_i:i\in J}\cap\gen{t_i:i\in K}
       =\gen{t_i:i\in J\cap K}\qquad(J,K\subseteq I).
\]
A string marking also satisfies $[t_i,t_j]=1$ when $|i-j|>1$.
We use the acquired recursive end criterion: a string of involutions
is a C-marking if its two strings obtained by deleting an end are
C-markings and their groups intersect in the group on their common
marks. For length at most two, distinct nonidentity involutions satisfy
the property. During the search each new involution lies outside the
group of all previously assigned marks, which ensures this base case.

\begin{lemma}\label{partial}
Let $I$ be the set of assigned labels, and suppose the string relations
hold on $I$. Write $I_1,\ldots,I_b$ for its maximal contiguous intervals,
$H_a=\gen{t_i:i\in I_a}$, and $H=\gen{t_i:i\in I}$. The partial marking
has the intersection property if and only if each interval marking has
that property and
\[
                         |H|=\prod_{a=1}^b |H_a|.                 \tag{1}
\]
\end{lemma}
\begin{proof}
Every pair of marks from different intervals has labels at distance at
least two. Thus the groups $H_a$ commute, and multiplication gives a
surjective homomorphism $\prod_a H_a\to H$. Equality (1) says exactly
that this finite surjection is injective. Under this identification,
the subgroup generated by any selected marks is the product of the
corresponding subgroups in its factors. Their intersections are
coordinatewise, so the interval intersection properties imply the full
partial property.

Conversely, restriction to an interval preserves the intersection
property. For each $a$, that property also gives
$H_a\cap\gen{t_i:i\in I\setminus I_a}=1$. If a product of one element
from each $H_a$ is the identity, its $a$th element belongs to this
intersection and hence is the identity. The multiplication map is
injective, proving (1).
\end{proof}

\begin{lemma}\label{candidates}
Fix a partial marking on $I$, with generated group $H$, and let $j$ be
the next label. Put
\[
 S=C_A(t_i:i\in I),\qquad
 Q=C_A(t_i:i\in I,\ |i-j|>1).
\]
For each $Q$-class of involutions contained in $G$, choose one member
$y$, and choose representatives $d$ for $C_Q(y)\backslash Q/S$.
The elements $y^d$ form exactly one representative of each $S$-orbit of
involutions in $G\cap Q$.
\end{lemma}
\begin{proof}
We have $S\leq Q$, and normality of $G$ makes membership in $G$
constant on each $Q$-class. Within the class of $y$, equality
$y^{d'}=(y^d)^s$ for some $s\in S$ holds precisely when
$d'\in C_Q(y)dS$. This is the acquired double coset argument, whose
hypotheses do not require the assigned labels to be an initial interval.
Different $Q$-classes cannot merge under the subgroup $S$.
\end{proof}

Discard a candidate $x$ if it lies in $H$, if it commutes with an
already assigned adjacent mark, or if the enlarged partial marking
fails Lemma~\ref{partial}. These tests are constant on its $S$-orbit.
The adjacent test is appropriate under our simplicity hypothesis. If a
full generating string of independent involutions had a commuting
adjacent pair, the string would split there into two nontrivial
commuting subgroups. Each would be normal in the generated group.
Simplicity would force each to equal $G$, forcing $G$ to be abelian,
a contradiction. For a general nonsimple group this discard must be
removed; it is not part of the general definition of a C-marking.

\begin{theorem}\label{enumeration}
Starting with the empty marking, apply Lemma~\ref{candidates} in any
fixed label order, retain exactly the candidates passing the stated
tests, and replace $S$ by $C_S(x)$ at each step. Discard a branch that
generates $G$ before all labels are assigned. At the final label retain
exactly the markings generating $G$. This procedure produces precisely
one representative of each ordered $\operatorname{Aut}(G)$-orbit of
rank-$r$ string C-markings of $G$.
\end{theorem}
\begin{proof}
The action induced by $S$ is the full stabilizer of the partial marking
under the assumed surjection from $A$. Lemma~\ref{candidates} therefore
enumerates the possible extensions of this orbit without repetition.
The string conditions are imposed by $Q$ and by the adjacent test;
Lemma~\ref{partial} imposes precisely the required intersection
property on every retained partial marking. A restriction of any full
C-marking has that property, and its next mark cannot belong to $H$:
otherwise the intersection property for its singleton and the old
label set would fail. Thus none of the discards loses a full marking.
In particular, once $H=G$, no further independent mark is possible.

The stabilizer of the extended marking is exactly $C_S(x)$, so
induction proves existence and uniqueness of the orbit representative
at every depth. At the last depth, requiring generation completes the
proof. The label order need not be changed during this induction.
\end{proof}

Reversal is an involution on the set of ordered automorphism orbits.
If there are $n$ ordered orbits and $s$ are equivalent to their
reversals, the number modulo duality is $(n+s)/2$. The program tests
this equivalence by exact conjugation of the full tuple in $A$.
The independent tuple checker also identifies every reversal partner
and verifies that its pairing is an involution.

\begin{lemma}\label{bound}
Let $P$ be a Sylow $2$-subgroup of $G$, and suppose the largest order of
an elementary abelian subgroup of $P$ is $2^e$. Every string C-marking
of $G$ has rank at most $2e$.
\end{lemma}
\begin{proof}
The odd-numbered marks commute. By the intersection property each adds
an independent involution, so they generate an elementary abelian
group of order $2^{\lceil r/2\rceil}$. This group lies in a conjugate
of $P$, proving $\lceil r/2\rceil\leq e$.
\end{proof}
The finite implementation transports $P$ by an actual isomorphism to a
polycyclic group, computes all its subgroup conjugacy classes, and
selects those of exponent at most two. A finite group of exponent two
is elementary abelian, so the maximum of their orders is exactly the
quantity in the lemma. Conjugacy preserves order and exponent. This
supplies a proved stopping rank from the computed subgroup list;
neither a source table nor a presumed odd-field rank theorem supplies it.

\paragraph{The prefix control.}
The acquired prefix enumeration is retained as a separate comparison.
Its present minimum-rank lookahead is also exact. At a prefix of length
$m\leq2$, a rank at least four requires a mark after the next mark.
That later mark centralizes the whole current prefix and lies outside
$H$. Let $K$ be the normal closure in this centralizer of all its
involution classes lying in $G$. If $K\leq H$, the extension is
impossible. Generating tuples of length below four need not be stored,
since generation already prevents independent extension. No predicted
census count enters either traversal.

\paragraph{Finite control and limits.}
For the identified group $G_2(3)$, the compatible automorphism action
has degree $1456$, $|G|=4245696$ and $|A|=8491392$. The Sylow subgroup
has order $64$ and its complete list has $84$ subgroup conjugacy
classes; the largest elementary abelian order is $8$. Thus ranks
$4,5,6$ exhaust the ranks at least four. Both traversals find exactly
two ordered rank-four orbits, paired by reversal, and no higher
orbits. Their actual tuples, not just their counts, are matched by
exact conjugation in $A$. Each tuple is separately checked for
generation, all string relations, and all $4^4=256$ subset
intersections. Its type is $[3,8,3]$; the two end groups have orders
$11232$ and $336$, in reverse orders for the dual tuple. Their named
isomorphism types are not certified in this methodological component.

Two separate computations for $G_2(5)$ were stopped at their recorded
120-second deadlines. They certify no census. The present packet
therefore proves the general traversal theorem under its stated
premises and records a finite $q=3$ control. It neither completes the
finite tables for F004 nor earns whole-original credit.

\paragraph{Attribution.}
The stabilizer enumeration is the method of Leemans--Mixer (2012),
Section~5, there attributed to Hartley--Hulpke (2010). We retain its
complete acquired proof and the standard end intersection criterion.
The contribution here is the justified change of label order and the
test for disconnected partial selections. The double coset and Sylow
arguments are standard finite group arguments. No claim of priority,
new open-conjecture solution, or formal Lean verification is made.
\end{document}
